\begingroup
\renewcommand{\thesection}{13C}
\section{Scalar Complexity and Relational Orientation}\label{sec:scalarmode}
\status{descriptive separation and conditional covariance; the native Mode theorem is unchanged}
The intended reading of the scalar is \emph{relational complexity of the thing}. A phase or orientation changes how that structure is presented, not necessarily how much structure it contains. In this sense the two descriptions are orthogonal:
\[
\boxed{\mathcal C(R)}\ \perp_{\rm desc}\
\boxed{\text{phase, Mode, and relational orientation}}.
\]
The subscript specifies descriptive separation, not a supplied inner product, statistical independence, or a proof that all values can be freely combined.

Suppose an admitted presentation action $g$ takes $R$ to an equivalent registered presentation. Definition~4A.2 then gives
\begin{equation}
\mathcal C(g\cdot R)=\mathcal C(R),\qquad S(g\cdot R)=S(R).
\tag{RC7}\label{eq:scalarinvariance}
\end{equation}
When the $S_3$ frame action of Section~\ref{sec:mode} has been attached to such objects on a common domain, a compatible descriptive chart $(c,p,m)$ has
\begin{equation}
\begin{aligned}
\Cevt:(c,p,m)&\longmapsto(c,p+1,m),\\
\Mevt:(c,p,m)&\longmapsto(c,-p,1-m),
\end{aligned}
\qquad p\in\Z_3,\quad m\in\Z_2.
\tag{RC8}\label{eq:complexitymodechart}
\end{equation}
The first coordinate is unchanged; the last two retain exactly the semidirect action (M4). This is a compatible extension under the stated presentation hypothesis, not a newly derived cross-register map.

A reflection can reverse a signed coordinate in a real representation without reversing the nonnegative complexity of its underlying object. In particular, neither $\mathcal C\mapsto-\mathcal C$ nor $S\mapsto-S$ follows from the retained Mode theorem. A signed amplitude, if required, must be separately defined and its covariance proved. Likewise a genuinely structure-changing operation need not preserve $\mathcal C$ merely because it is informally called a reversal.

The motivating sketch placed scalar complexity beside an oriented return about a distinguished null. It helps keep the descriptive roles separate, but no identification of that null with the native phase center is used in a definition or proof here. The null-well, event-frame, and scalar domains require an explicit common-domain construction before the picture can become a native theorem.

\subsection*{When orthogonality becomes geometric}
Descriptive separation does not supply a metric. If an independently declared orientation representation has $\mathbf n\cdot\mathbf n=1$, differentiating gives the actual metric statement $\mathbf n\cdot\dot{\mathbf n}=0$. Section~\ref{sec:gyroscopic} proves the resulting tangent and stabilizer geometry, and states the additional conditions for uniform precession. This metric example neither redefines the capitalized native Mode nor proves that a pure mechanical orientation change preserves an as-yet unconstructed $\mathcal C$.
\endgroup
\setcounter{section}{13}
